% Figure: a resonance peak with the half-power bandwidth marked.
% Magnitude of a second-order bandpass |H| against normalised frequency,
% showing the peak, the -3 dB (half-power) points, and the bandwidth that
% sets the quality factor Q = f_0 / Delta_f.
\begin{figure}[htbp]
\centering
\resizebox{0.82\linewidth}{!}{%
\begin{tikzpicture}
\begin{axis}[
  width=12cm, height=7.4cm,
  xlabel={frequency $f/f_0$},
  ylabel={magnitude $|H|$ (relative to peak)},
  xmin=0.6, xmax=1.4, ymin=0, ymax=1.12,
  xtick={0.7, 0.8, 0.9, 1.0, 1.1, 1.2, 1.3},
  ytick={0, 0.25, 0.5, 0.707, 1.0},
  yticklabels={$0$, $0.25$, $0.5$, $0.707$, $1$},
  axis lines=left,
  tick align=outside,
  every axis plot/.append style={very thick},
  label style={font=\footnotesize},
  tick label style={font=\footnotesize},
  clip=false,
]
  % Second-order bandpass magnitude, normalised to peak 1:
  % |H(u)| = 1 / sqrt( 1 + Q^2 (u - 1/u)^2 ), u = f/f0, Q = 10.
  % At u away from 1 the argument stays positive, so sqrt is safe.
  \addplot[engblue, domain=0.6:1.4, samples=300] {1/sqrt(1 + 100*(x - 1/x)^2)};
  % half-power line
  \addplot[engblack!55, dashed, domain=0.6:1.4, samples=2] {0.707};
  % peak marker
  \addplot[engred, only marks, mark=*, mark size=1.6pt] coordinates {(1.0, 1.0)};
  % half-power points: Q=10 -> Delta(f/f0) ~ 1/Q = 0.1, at ~0.951 and ~1.051
  \addplot[engred, only marks, mark=*, mark size=1.6pt]
    coordinates {(0.9512, 0.707) (1.0512, 0.707)};
  \draw[<->, engamber, thick]
    (axis cs:0.9512, 0.63) -- (axis cs:1.0512, 0.63);
  \node[font=\footnotesize, text=engamber, anchor=north]
    at (axis cs:1.0, 0.60) {bandwidth $\Delta f$};
  \node[font=\footnotesize, text=engblack!70, anchor=west]
    at (axis cs:1.16, 0.707) {half-power ($-3\,\si{\deci\bel}$)};
  \node[font=\footnotesize, text=engred, anchor=south]
    at (axis cs:1.0, 1.0) {peak at $f_0$};
\end{axis}
\end{tikzpicture}%
}%
\caption[Resonance peak and half-power bandwidth]{The magnitude of a
second-order resonant response against frequency, normalised so the
peak is one. The two frequencies at which the magnitude falls to
$1/\sqrt{2} \approx 0.707$ of the peak, the half-power or
$-3\,\si{\deci\bel}$ points, bracket the bandwidth $\Delta f$. The
quality factor is the ratio $Q = f_0 / \Delta f$, so a narrow peak (a
small $\Delta f$) is a high-$Q$ resonance and a broad peak a low-$Q$
one. Reading the peak frequency and the two half-power crossings off
the curve fixes both $f_0$ and $Q$ without a nonlinear fit, the same
two-reading discipline the load-cell and moving-coil case studies
used on their own curves.}
\label{fig:vol02:ch02:resonance-q}
\end{figure}
